% LIGHT system map for the basic network: energy flow through the Kr/F2
% kinetics with KrF in its B and C states. Every number comes from
% system_map_basic_numbers.tex, which make_system_map_basic.py writes from
% the model. Build from LIGHT/docs:
%   python3 make_system_map_basic.py && latexmk -pdf system_map_basic.tex
% The v0 map (system_map.tex) draws the notes' list as written.
\documentclass[11pt]{article}
\usepackage[a4paper,landscape,margin=0.9cm,top=0.7cm,bottom=0.6cm]{geometry}
\usepackage[T1]{fontenc}
\usepackage[scaled=0.95]{helvet}
\renewcommand{\familydefault}{\sfdefault}
\usepackage{amsmath}
\usepackage{sansmath}
\usepackage{booktabs,tabularx,array}
\usepackage{xcolor}
\usepackage{tikz}
\usetikzlibrary{arrows.meta,calc,positioning}
\pagestyle{empty}
\setlength{\parindent}{0pt}

% LIGHT basic-network map numbers. GENERATED by make_system_map_basic.py
% from the model (0.1.0). Do not edit. Rerun the script instead.
\providecommand\lmdef[2]{\expandafter\def\csname lm@#1\endcsname{#2}}
\lmdef{e.P}{100.0}
\lmdef{e.pen}{11.2}
\lmdef{e.G1}{90.2}
\lmdef{e.C1}{11.7}
\lmdef{e.C2}{0.3}
\lmdef{e.G6}{8.9}
\lmdef{e.dU}{0.0}
\lmdef{e.cap}{0.0}
\lmdef{C1.kr}{10.0}
\lmdef{C1.heat}{1.7}
\lmdef{krs.in}{100.2}
\lmdef{krs.harp}{61.8}
\lmdef{krs.G2}{37.9}
\lmdef{krs.left}{0.5}
\lmdef{brs.harp}{61.6}
\lmdef{brs.G2}{37.9}
\lmdef{harp.B}{21.4}
\lmdef{harp.C}{13.1}
\lmdef{harp.F}{5.5}
\lmdef{harp.heat}{21.8}
\lmdef{G2.ion}{26.8}
\lmdef{C2.ion}{0.3}
\lmdef{mix.net}{2.9}
\lmdef{mix.netC}{2.9}
\lmdef{mix.heat}{0.0}
\lmdef{mix.BC}{11.5}
\lmdef{mix.CB}{14.4}
\lmdef{B.G4}{6.4}
\lmdef{B.G5}{13.8}
\lmdef{B.Q}{4.0}
\lmdef{B.QF2}{1.0}
\lmdef{B.QKR}{0.0}
\lmdef{B.QE}{3.0}
\lmdef{B.left}{0.0}
\lmdef{C.G4C}{1.1}
\lmdef{C.G5C}{3.6}
\lmdef{C.Q}{5.5}
\lmdef{C.QF2}{0.8}
\lmdef{C.QKR}{1.1}
\lmdef{C.QE}{3.6}
\lmdef{C.left}{0.0}
\lmdef{G4.hv}{5.2}
\lmdef{G4.F}{0.8}
\lmdef{G4.heat}{0.4}
\lmdef{G4C.hv}{0.8}
\lmdef{G4C.F}{0.1}
\lmdef{G4C.heat}{0.1}
\lmdef{QB.F}{0.8}
\lmdef{QB.heat}{3.2}
\lmdef{QC.F}{0.9}
\lmdef{QC.heat}{4.6}
\lmdef{tri.k2f}{11.6}
\lmdef{tri.heat}{5.8}
\lmdef{k2f.C3}{5.6}
\lmdef{k2f.Q2F}{6.0}
\lmdef{k2f.left}{0.0}
\lmdef{brk.C3}{48.0}
\lmdef{brk.Q2F}{51.9}
\lmdef{C3.hv}{4.1}
\lmdef{C3.F}{1.1}
\lmdef{C3.heat}{0.4}
\lmdef{Q2F.F}{1.2}
\lmdef{Q2F.heat}{4.8}
\lmdef{out.248}{5.2}
\lmdef{out.275}{0.8}
\lmdef{out.415}{4.1}
\lmdef{out.heat}{42.9}
\lmdef{out.elastic}{8.9}
\lmdef{out.ions}{27.1}
\lmdef{out.F}{10.4}
\lmdef{out.other}{0.5}
\lmdef{version}{0.1.0}
\lmdef{p}{1}
\lmdef{xF2}{0.5}
\lmdef{Tg0}{300}
\lmdef{P}{300}
\lmdef{pulse}{1}
\lmdef{tend}{2}
\lmdef{ne}{10^{15}}
\lmdef{N}{2.45{\times}10^{19}}
\lmdef{Edep}{0.30}
\lmdef{Te}{2.03}
\lmdef{TeMid}{2.00}
\lmdef{rTail}{0.30}
\lmdef{TtMid}{0.60}
\lmdef{TgEnd}{604}
\lmdef{pEnd}{2.03}
\lmdef{F2burn}{715}
\lmdef{F2left}{0.04}
\lmdef{krsFormed}{1.89{\times}10^{17}}
\lmdef{krfFormed}{1.05{\times}10^{17}}
\lmdef{k2fFormed}{5.30{\times}10^{16}}
\lmdef{krf.eff}{28.0}
\lmdef{brf.G4}{18.7}
\lmdef{brf.G4C}{3.2}
\lmdef{brf.tri}{50.5}
\lmdef{brf.e}{19.0}
\lmdef{brf.F2}{5.4}
\lmdef{brf.Kr}{3.2}
\lmdef{brf.QB}{11.6}
\lmdef{brf.QC}{16.1}
\lmdef{k.harp}{8.47{\times}10^{-10}}
\lmdef{pc.harpY}{90}
\lmdef{pc.harpB}{56}
\lmdef{pc.harpC}{34}
\lmdef{pc.harpX}{10}
\lmdef{pc.fC}{38}
\lmdef{k.G2}{5.0{\times}10^{-10}}
\lmdef{tauB}{6.99}
\lmdef{tauC}{50}
\lmdef{tauC3}{220}
\lmdef{gapBC}{80}
\lmdef{k.BC}{1.05{\times}10^{-11}}
\lmdef{Keq}{1.10}
\lmdef{k.G5}{1.9{\times}10^{-30}}
\lmdef{k.G5C}{3.79{\times}10^{-31}}
\lmdef{tpow.G5}{-3.32}
\lmdef{k.Q2F}{1.2{\times}10^{-10}}
\lmdef{hv.G4C}{4.5}
\lmdef{hv.G4}{5.00}
\lmdef{notes.krs}{101.9}
\lmdef{notes.krp}{27.1}
\lmdef{notes.emit}{10.1}
\lmdef{om.kr2}{8.0}
\lmdef{om.ions}{5.3}
\lmdef{om.electron}{3.8}
\lmdef{om.f2e}{5.7}
\lmdef{om.all}{8.0}
\lmdef{om.basic}{5.2}
\lmdef{om.v0}{16.8}

\makeatletter
\newcommand\lm[1]{\ifcsname lm@#1\endcsname\csname lm@#1\endcsname
  \else\lm@missing{#1}\fi}
\newcommand\lm@missing[1]{\errmessage{system_map_basic_numbers.tex has no
  number '#1'. Rerun make_system_map_basic.py}}
\makeatother

% categorical colors, checked with the dataviz validator (all pairs, white
% surface): blue, yellow, aqua, green, violet, red. KrF(C) is green beside
% KrF(B)'s aqua, the same molecule in its other state (normal-vision
% Delta E 15.6). Ions and F atoms are neutral greys (stored energy). Every
% band and tag also carries text.
\definecolor{cE}{HTML}{2A78D6}   % electron gas
\definecolor{cX}{HTML}{EDA100}   % Kr* excitation
\definecolor{cB}{HTML}{1BAF7A}   % KrF(B) and 248 nm
\definecolor{cC}{HTML}{008300}   % KrF(C) and the C-A band
\definecolor{cT}{HTML}{4A3AA7}   % Kr2F* and 415 nm
\definecolor{cH}{HTML}{E34948}   % heat
\definecolor{cI}{HTML}{52514E}   % ions (stored)
\definecolor{cF}{HTML}{A8A69E}   % F atoms (stored)
\definecolor{ink}{HTML}{0B0B0B}
\definecolor{ink2}{HTML}{52514E}
\definecolor{muted}{HTML}{898781}

% chemistry and units in running text or math
\newcommand{\unit}[1]{\ensuremath{\,\mathrm{#1}}}
\newcommand{\el}{\ensuremath{\mathrm{e}^-}}
\newcommand{\Krs}{Kr\ensuremath{^*}}
\newcommand{\Krp}{Kr\ensuremath{^+}}
\newcommand{\KtF}{Kr\ensuremath{_2}F\ensuremath{^*}}
\newcommand{\Ft}{F\ensuremath{_2}}
\newcommand{\ra}{\ensuremath{\to}\ }
\newcommand{\hnu}{\ensuremath{h\nu}}
\newcommand{\Edep}{\ensuremath{E_\mathrm{dep}}}
\newcommand{\pc}[1]{\lm{#1}\,\%}
\newcommand{\cm}[1]{\ensuremath{\,\mathrm{cm^{#1}s^{-1}}}}

\tikzset{
  band/.style={line width={max(0.6pt,(#1)*0.03mm)},
    -{Triangle[width=2.2pt 1.8,length=2pt 1.0]}},
  thin arrow/.style={line width=0.6pt,
    -{Triangle[width=2.2pt 1.8,length=2pt 1.0]}},
  blab/.style={font=\scriptsize, text=ink, inner sep=0.8pt, fill=white,
    fill opacity=0.85, text opacity=1},
  nohyph/.style={execute at begin node={\hyphenpenalty=10000
    \exhyphenpenalty=10000}},
  proc/.style={nohyph, draw=ink2!55, fill=white, rounded corners=2pt,
    text width=2.6cm, align=left, font=\scriptsize, inner sep=3pt,
    text=ink},
  spec/.style={nohyph, draw=#1, fill=#1!12, rounded corners=5pt,
    align=center, font=\scriptsize, text=ink, inner sep=3pt,
    text width=2.05cm, line width=0.8pt},
  inbox/.style={nohyph, draw=ink2!55, fill=white, rounded corners=2pt,
    text width=2.5cm, align=left, font=\scriptsize, inner sep=4pt,
    text=ink},
  tag/.style={draw=#1, fill=#1!14, rounded corners=1.5pt,
    font=\tiny, text=ink, inner sep=1.6pt},
  hdr/.style={font=\footnotesize\bfseries, text=ink2},
  note/.style={nohyph, font=\tiny, text=ink2},
}

\begin{document}
\sansmath

{\large\bfseries LIGHT system map, basic network: energy flow through the
Kr/F$_2$ kinetics}\par
\vspace{2pt}
{\footnotesize Reference point: \lm{p}\unit{atm} of Kr with \pc{xF2} F$_2$ at
\lm{Tg0}\unit{K}, $P=\lm{P}\unit{kW\,cm^{-3}}$ for \lm{pulse}\unit{\mu s}
so that $\Edep=\lm{Edep}\unit{J\,cm^{-3}}$, $n_e=\lm{ne}\unit{cm^{-3}}$
held, run to \lm{tend}\unit{\mu s}. Band width is proportional to the
energy flow. Numbers are percent of \Edep, rounded to 0.1.}\par
\vspace{2pt}

\begin{tikzpicture}[x=1cm,y=1cm]
% ================================================================ frames
\node[hdr] at (1.4,15.7) {INPUT};
\node[hdr] at (26.3,15.7) {OUTPUT at $t=\lm{tend}\unit{\mu s}$};
\node[note] at (26.3,15.32) {percent of \Edep, sums to 100};
\draw[draw=ink2!35, fill=black!2.5, rounded corners=6pt]
  (3.1,0.25) rectangle (24.25,15.45);
\node[hdr, anchor=west] at (3.3,15.08)
  {UNIFORM INITIAL GAS\quad 0-D, constant volume};

% ================================================================ inputs
\node[inbox, line width=0.9pt, draw=cE] (pow) at (1.4,12.5) {%
  \textbf{Power in}\\
  $P=\lm{P}\unit{kW\,cm^{-3}}$\\ for \lm{pulse}\unit{\mu s}, uniform,\\
  into the electrons\\
  $\Edep=\lm{Edep}\unit{J\,cm^{-3}}$\\
  \textcolor{ink2}{range 100 to 500}};
\node[inbox, dashed] (em) at (1.4,9.7) {%
  \textbf{or an EM wave}\\
  $E_\mathrm{rms}$ at frequency $f$\\
  absorbed by electrons\\
  \textcolor{ink2}{not used at this point}};
\node[inbox] (inj) at (1.4,7.1) {%
  \textbf{$n_e$ injection}\\
  $n_e=\lm{ne}\unit{cm^{-3}}$, held\\
  \textcolor{ink2}{range $10^{14}$ to $10^{16}$}};
\node[inbox] (fill) at (1.4,4.4) {%
  \textbf{Initial gas fill}\\
  Kr $+$ \pc{xF2} F$_2$\\ \lm{p}\unit{atm}, \lm{Tg0}\unit{K}\\
  \mbox{$N=\lm{N}\unit{cm^{-3}}$}\\
  \textcolor{ink2}{range 0.1 to 2\unit{atm}}};

% ================================================================ electron gas
\node[draw=cE, fill=cE!10, rounded corners=5pt, line width=0.9pt,
  minimum width=1.6cm, minimum height=11.1cm, anchor=north]
  (eg) at (4.0,14.6) {};
\node[align=center, font=\scriptsize, text=ink, text width=1.5cm,
  anchor=north] at ([yshift=-4pt]eg.north) {\textbf{electron}\\
  \textbf{gas}\\[4pt] Maxwellian\\ bulk\\ $T_e$ \lm{Te}\unit{eV}\\
  \textcolor{ink2}{pulse mean}\\[6pt] tail above\\ 9.915\unit{eV}\\
  $T_t$ \lm{TtMid}\unit{eV}\\ \textcolor{ink2}{mid-pulse}\\[6pt]
  $U_e=\tfrac32 n_e T_e$};
\node[align=center, note, text width=1.3cm, anchor=south]
  at ([yshift=3pt]eg.south) {in: power \lm{e.P}\\ and Penning
  \lm{e.pen}};

% ================================================================ e-impact rows
\node[note, anchor=south west, text width=5.3cm, align=left]
  at (5.45,14.05) {$Q=k\,n_e[\mathrm{Kr}]$, $k=\langle\sigma v\rangle$ over
  the two-temperature EEDF, Biagi v7.1 cross sections};
\node[proc, text width=2.4cm] (G1) at (6.75,13.3) {%
  \textbf{G1}\quad \el{} + Kr\\ \ra \el{} + \Krs(5s)\\
  \textcolor{ink2}{threshold 9.915\unit{eV}}};
\node[proc, text width=2.4cm] (C1) at (6.75,11.05) {%
  \textbf{C1}\quad \el{} + Kr\\ \ra Kr(5p, 4d, \dots)\\
  \textcolor{ink2}{cascades to \Krs(5s)}};
\node[proc, text width=2.4cm] (C2) at (6.75,5.6) {%
  \textbf{C2}\quad \el{} + Kr\\ \ra \Krp{} + 2\el\\
  \textcolor{ink2}{threshold 13.996\unit{eV}}};
\node[proc, text width=2.4cm] (G6) at (6.75,4.1) {%
  \textbf{G6}\quad elastic\\ \el{} + Kr\\
  \textcolor{ink2}{$(2m/M)(1-T_g/T_e)$}};

% ================================================================ Kr*
\node[spec=cX] (KS) at (9.75,12.2) {%
  \textbf{\Krs(5s)}\\ formed\\ $\lm{krsFormed}\unit{cm^{-3}}$};

% ================================================================ harpoon, G2
\node[proc] (H) at (12.9,13.1) {%
  \textbf{G3, G3C, G3X}\\ \Krs{} + \Ft\\
  \textcolor{ink2}{$\lm{k.harp}\cm{3}$}\\
  \ra KrF(B) + F, \pc{pc.harpB}\\
  \ra KrF(C) + F, \pc{pc.harpC}\\
  \ra Kr + 2F, \pc{pc.harpX}\\
  takes \pc{brs.harp} of \Krs};
\node[proc] (G2) at (12.9,9.85) {%
  \textbf{G2}\quad \Krs{} + \Krs\\ \ra \Krp{} + Kr + \el\\
  \textcolor{ink2}{$\lm{k.G2}\cm{3}$}\\
  takes \pc{brs.G2} of \Krs};

% ================================================================ KrF(B), KrF(C), Kr+
\node[spec=cB] (KB) at (16.05,13.45) {%
  \textbf{KrF(B)}\\ upper laser\\ state};
\node[spec=cC] (KC) at (16.05,9.75) {%
  \textbf{KrF(C)}\\ \lm{gapBC}\unit{cm^{-1}}\\ below B};
\node[spec=cI] (KI) at (16.05,5.6) {%
  \textbf{\Krp{} ions}\\ {\tiny stored, not recombined here}};

% ================================================================ B and C fates
\node[proc] (G4) at (19.2,14.45) {%
  \textbf{G4}\quad KrF(B)\\ \ra Kr + F + \hnu\\
  \textcolor{ink2}{248\unit{nm}, $\tau_B=\lm{tauB}\unit{ns}$}};
\node[proc] (QB) at (19.2,12.85) {%
  \textbf{QE, QF2, QKR}\\ quench KrF(B) by\\
  \el\,\lm{B.QE}, \Ft\,\lm{B.QF2}, Kr\,\lm{B.QKR}};
\node[proc] (TR) at (19.2,11.2) {%
  \textbf{G5, G5C}\quad KrF(B, C)\\ + 2Kr \ra \KtF{} + Kr\\
  \textcolor{ink2}{$\times\,(T_g/300)^{\lm{tpow.G5}}$}};
\node[proc] (QC) at (19.2,9.6) {%
  \textbf{QEC, QF2C, QKRC}\\ quench KrF(C) by\\
  \el\,\lm{C.QE}, \Ft\,\lm{C.QF2}, Kr\,\lm{C.QKR}};
\node[proc] (G4C) at (19.2,7.95) {%
  \textbf{G4C}\quad KrF(C)\\ \ra Kr + F + \hnu\\
  \textcolor{ink2}{$\approx$275\unit{nm}, $\tau_C=\lm{tauC}\unit{ns}$}};

% ================================================================ Kr2F*
\node[proc] (C3) at (22.7,13.0) {%
  \textbf{C3}\quad \KtF\\ \ra 2Kr + F + \hnu\\
  \textcolor{ink2}{$\approx$415\unit{nm}, $\tau=\lm{tauC3}\unit{ns}$}\\
  \pc{brk.C3} of \KtF};
\node[spec=cT] (K2) at (22.7,11.2) {%
  \textbf{\KtF}\\ {\unsansmath$4^2\Gamma$} trimer\\
  $\lm{k2fFormed}\unit{cm^{-3}}$};
\node[proc] (Q2) at (22.7,9.4) {%
  \textbf{Q2F}\quad \KtF{} + \Ft\\
  \textcolor{ink2}{$\lm{k.Q2F}\cm{3}$}\\
  \pc{brk.Q2F} of \KtF};

% ================================================================ bands
% inputs into the gas
\draw[band=\lm{e.P}, cE] (pow.east) -- (pow.east -| eg.west)
  node[blab, midway, above=2.2pt] {\lm{e.P}};
\draw[thin arrow, cE, dashed] (em.east) -- (em.east -| eg.west);
\draw[thin arrow, ink2] (inj.east) -- (inj.east -| eg.west)
  node[blab, midway, above=1pt, font=\tiny] {sets $n_e$};
\draw[thin arrow, ink2] (fill.east) -- (fill.east -| 3.1,0)
  node[blab, midway, above=1pt, font=\tiny] {fills};
% electron gas out to the four electron-impact rows
\foreach \n/\k in {G1/e.G1, C1/e.C1, C2/e.C2, G6/e.G6}{
  \draw[band=\lm{\k}, cE] (\n.west -| eg.east) -- (\n.west)
    node[blab, midway, above=1.5pt] {\lm{\k}};
}
% into Kr*
\draw[band=\lm{e.G1}, cX] (G1.east) -- (KS.north west)
  node[blab, midway, above right=0pt and 0pt] {\lm{e.G1}};
\draw[band=\lm{C1.kr}, cX] (C1.east) -- (KS.south west)
  node[blab, midway, below right=0pt and 0pt] {\lm{C1.kr}};
% out of Kr*
\draw[band=\lm{krs.harp}, cX] (KS.north east) -- (H.west |- KS.north east)
  node[blab, midway, above=2pt] {\lm{krs.harp}};
\draw[band=\lm{krs.G2}, cX] (KS.south east) -- (G2.west)
  node[blab, midway, below left=0pt and 0pt] {\lm{krs.G2}};
% the harpoon into KrF(B) and KrF(C)
\draw[band=\lm{harp.B}, cB] (H.east |- KB.west) -- (KB.west)
  node[blab, midway, above=2pt] {\lm{harp.B}};
\draw[band=\lm{harp.C}, cC] (H.south east) -- (KC.north west)
  node[blab, pos=0.5, left=2pt] {\lm{harp.C}};
% G2 into the ions and back to the electrons
\draw[band=\lm{G2.ion}, cI] (G2.south east) -- (KI.north west)
  node[blab, pos=0.45, right=3pt] {\lm{G2.ion}};
\draw[band=\lm{e.pen}, cE, rounded corners=5pt]
  (G2.south west) -- (G2.south west |- 0,8.55) -- (eg.east |- 0,8.55)
  node[blab, pos=0.5, below=2pt] {Penning returns \lm{e.pen} to the
  electrons, 5.83\unit{eV} per event};
\draw[band=\lm{C2.ion}, cI] (C2.east) -- (KI.west)
  node[blab, pos=0.5, above=1.5pt] {\lm{C2.ion}};
% B-C mixing, drawn as its net flow
\draw[band=\lm{mix.net}, cC] (KC.north) -- (KB.south);
\node[note, anchor=west, align=left, text width=1.1cm]
  at (16.15,11.3) {\textbf{BC, CB}\\ mixing by\\ Kr, net\\ C \ra B
  \textcolor{ink}{\lm{mix.net}}};
% out of KrF(B)
\draw[band=\lm{B.G4}, cB] (KB.north east) -- (G4.west)
  node[blab, midway, above left=0pt and 0pt] {\lm{B.G4}};
\draw[band=\lm{B.Q}, cB] (KB.east) -- (QB.west)
  node[blab, pos=0.5, above=1pt] {\lm{B.Q}};
\draw[band=\lm{B.G5}, cB] (KB.south east) -- ([yshift=5pt]TR.west)
  node[blab, pos=0.42, left=1pt] {\lm{B.G5}};
% out of KrF(C)
\draw[band=\lm{C.G5C}, cC] (KC.north east) -- ([yshift=-5pt]TR.west)
  node[blab, pos=0.42, below right=0pt and 0pt] {\lm{C.G5C}};
\draw[band=\lm{C.Q}, cC] (KC.east) -- (QC.west)
  node[blab, pos=0.4, below=1pt] {\lm{C.Q}};
\draw[band=\lm{C.G4C}, cC] (KC.south east) -- (G4C.west)
  node[blab, midway, below left=0pt and 0pt] {\lm{C.G4C}};
% trimer and Kr2F*
\draw[band=\lm{tri.k2f}, cT] (TR.east) -- (K2.west)
  node[blab, midway, above=2pt] {\lm{tri.k2f}};
\draw[band=\lm{k2f.C3}, cT] (K2.north) -- (C3.south)
  node[blab, midway, right=3pt] {\lm{k2f.C3}};
\draw[band=\lm{k2f.Q2F}, cT] (K2.south) -- (Q2.north)
  node[blab, midway, left=3pt] {\lm{k2f.Q2F}};
% light and ions to the outputs
\draw[band=\lm{G4.hv}, cB] (G4.east) -- (G4.east -| 24.7,0)
  node[blab, pos=0.75, above=2pt] {\lm{G4.hv}};
\draw[band=\lm{C3.hv}, cT] (C3.east) -- (C3.east -| 24.7,0)
  node[blab, midway, above=1.5pt] {\lm{C3.hv}};
\draw[band=\lm{G4C.hv}, cC] (G4C.east) -- (G4C.east -| 24.7,0)
  node[blab, pos=0.75, above=1.5pt] {\lm{G4C.hv}};
\draw[band=\lm{out.ions}, cI] (KI.east) -- (KI.east -| 24.7,0)
  node[blab, pos=0.8, above=2.2pt] {\lm{out.ions}};

% ================================================================ loss tags
\node[tag=cH, anchor=north east] at ([yshift=-2pt]C1.south east)
  {heat \lm{C1.heat}};
\node[tag=cH, anchor=north east] at ([yshift=-2pt]G6.south east)
  {heat \lm{e.G6}, elastic};
\node[tag=cF, anchor=north west] at ([yshift=-2pt]H.south west)
  {F atoms \lm{harp.F}};
\node[tag=cH, anchor=north east] at ([yshift=-2pt]H.south east)
  {heat \lm{harp.heat}};
\foreach \n/\f/\h in {G4/G4.F/G4.heat, QB/QB.F/QB.heat, QC/QC.F/QC.heat,
    G4C/G4C.F/G4C.heat, C3/C3.F/C3.heat, Q2/Q2F.F/Q2F.heat}{
  \node[tag=cF, anchor=north west] at ([yshift=-2pt]\n.south west)
    {F atoms \lm{\f}};
  \node[tag=cH, anchor=north east] at ([yshift=-2pt]\n.south east)
    {heat \lm{\h}};
}
\node[tag=cH, anchor=north east] at ([yshift=-2pt]TR.south east)
  {heat \lm{tri.heat}};

% ================================================================ groups
\node[draw=ink2!60, dashed, rounded corners=3pt, fill=white,
  align=left, font=\scriptsize, text=ink, anchor=south west,
  text width=9.3cm, inner sep=5pt] at (5.45,0.45) {%
  \textbf{Kr/F$_2$ channels not in the basic network.} The 248-nm share
  of \Edep{} with one optional group switched on. Measured with
  \texttt{run\_v0.py -{}-omissions}.\\[3pt]
  \begin{tabular}{@{}l@{\hspace{3pt}}r@{\hspace{14pt}}l@{\hspace{3pt}}r
    @{\hspace{14pt}}l@{\hspace{3pt}}r@{}}
  basic as built & \pc{om.basic} & \Krs{} + 2Kr \ra Kr$_2^*$ & \pc{om.kr2}
    & ion recombination & \pc{om.ions}\\
  \el{} on \Krs & \pc{om.electron} & \el{} on \Ft & \pc{om.f2e} &
    \textbf{all groups on} & \textbf{\pc{om.all}}\\
  \end{tabular}\\[3pt]
  The notes' six reactions as written (v0) give \pc{om.v0}.};

% ================================================================ legend
\begin{scope}[shift={(15.55,0.5)}, font=\scriptsize, text=ink]
  \node[anchor=west, font=\scriptsize\bfseries, text=ink2] at (0,3.3)
    {Band color is where the energy sits};
  \foreach \c/\t [count=\i] in {cE/electron gas, cX/\Krs{} excitation,
      cB/KrF(B) and 248-nm light, cC/KrF(C) and C-A light,
      cT/\KtF{} and 415-nm light, cI/\Krp{} ions}{
    \draw[\c, line width=1.3mm] (0,3.3-0.47*\i) -- ++(0.55,0);
    \node[anchor=west] at (0.62,3.3-0.47*\i) {\t};
  }
  \node[tag=cH, anchor=west] at (4.45,2.83) {heat};
  \node[anchor=west] at (5.15,2.83) {turns into heat};
  \node[tag=cF, anchor=west] at (4.45,2.36) {F atoms};
  \node[anchor=west] at (5.6,2.36) {left in free F};
  \draw[ink2, line width={10*0.03mm}] (4.45,1.89) -- ++(0.55,0);
  \node[anchor=west] at (5.05,1.89) {10\,\% band};
  \draw[ink2, line width={50*0.03mm}] (4.45,1.42) -- ++(0.55,0);
  \node[anchor=west] at (5.05,1.42) {50\,\% band};
\end{scope}

% ================================================================ ledger
\begin{scope}[font=\scriptsize, text=ink]
  \newcommand{\ledger}[5]{% y, color, label, key, bar options
    \node[anchor=south west, inner sep=1pt] at (24.75,#1) {#3};
    \node[anchor=south east, inner sep=1pt, font=\scriptsize\bfseries]
      at (27.85,#1) {\lm{#4}};
    \draw[ink2!30] (24.8,#1-0.1) -- (27.8,#1-0.1);
    \fill[#2, #5] (24.8,#1-0.1) rectangle ++({3.0*(\lm{#4})/100},-0.14);
  }
  \ledger{14.45}{cB}{248-nm emission}{out.248}{}
  \ledger{13.0}{cT}{415-nm emission}{out.415}{}
  \ledger{7.95}{cC}{C-A emission}{out.275}{}
  \ledger{5.6}{cI}{stored in \Krp{} ions}{out.ions}{}
  \ledger{4.6}{cH}{reaction heat}{out.heat}{}
  \ledger{3.7}{cH}{elastic losses}{out.elastic}{opacity=0.55}
  \ledger{2.8}{cF}{stored in F atoms}{out.F}{}
  \ledger{1.9}{muted}{other stores}{out.other}{}
  \node[note, anchor=north west, text width=3.1cm, align=left]
    at (24.75,1.45) {The heat and F tags inside the gas add up to these
    totals. Ledger closure is below $10^{-12}$ of \Edep.};
\end{scope}
\end{tikzpicture}

\vspace{1pt}
{\footnotesize The notes' outputs as flows, which overlap: to \Krs{}
\pc{notes.krs}, to \Krp{} \pc{notes.krp}, to KrF(B) and KrF(C)
\pc{krf.eff}, heat \pc{out.heat}, elastic losses \pc{out.elastic},
emission \pc{notes.emit}. Of the KrF formed,
$\lm{krfFormed}\unit{cm^{-3}}$, \pc{brf.G4} radiates at 248\unit{nm} and
\pc{brf.G4C} in the C-A band, \pc{brf.tri} forms the trimer, and
electrons, F$_2$ and Kr quench \pc{brf.e}, \pc{brf.F2} and \pc{brf.Kr}.
B-C mixing carries \pc{mix.BC} of \Edep{} from B to C and \pc{mix.CB}
back. At the end of the run $T_g=\lm{TgEnd}\unit{K}$ and
$p=\lm{pEnd}\unit{atm}$, and F$_2$ fell below 10\,\% of the fill at
\lm{F2burn}\unit{ns}.}\par
{\scriptsize\color{ink2} Generated from LIGHT v0 (\lm{version}), basic
network, by \texttt{docs/make\_system\_map\_basic.py}. No number on this
page is typed by hand. The rates and their sources are in
\texttt{accuracy\_overview.pdf}, Appendix A.}

\end{document}
